\documentclass[11pt]{amsart}

\usepackage[T1]{fontenc}
\usepackage[utf8]{inputenc}
\usepackage{lmodern}
\usepackage[margin=1.15in]{geometry}
\usepackage{amsmath,amssymb,amsthm,mathtools}
\usepackage{booktabs}
\usepackage[colorlinks=true,linkcolor=blue!50!black,citecolor=green!35!black,urlcolor=blue!50!black]{hyperref}
\usepackage{xcolor}

\newtheorem{theorem}{Theorem}
\newtheorem{corollary}[theorem]{Corollary}
\newtheorem{proposition}[theorem]{Proposition}
\newtheorem{lemma}[theorem]{Lemma}
\theoremstyle{definition}
\newtheorem{definition}[theorem]{Definition}
\theoremstyle{remark}
\newtheorem{remark}[theorem]{Remark}

\setlength{\emergencystretch}{3em}
\newcommand{\Z}{\mathbb{Z}}
\newcommand{\N}{\mathbb{N}}
\newcommand{\R}{\mathbb{R}}
\newcommand{\C}{\mathbb{C}}
\newcommand{\Gq}{G_q}
\newcommand{\bias}{\operatorname{bias}}
% Lean identifiers: typewriter, breakable after `.` `_` `/`
\def\UrlBreaks{\do\.\do\_\do\-\do\/\do\A\do\B\do\C\do\D\do\E\do\F\do\G\do\H\do\I\do\J\do\K\do\L\do\M\do\N\do\O\do\P\do\Q\do\R\do\S\do\T\do\U\do\V\do\W\do\X\do\Y\do\Z}
\newcommand{\lean}{\nolinkurl}
\newcommand{\abs}[1]{\left\lvert #1 \right\rvert}
\newcommand{\norm}[1]{\left\lVert #1 \right\rVert}

\title[Explicit, verified spectral bounds for prime Cayley graphs]{Explicit and formally verified spectral bounds\\ for prime Cayley graphs on $(\Z/q\Z)^\times$}

\author{[Author names to be added]}

\date{Draft of 27 November 2026}

\begin{document}

\begin{abstract}
Let $q\ge 1$ and let $\mu$ be a probability measure on integers coprime to $q$. Its bias
$\max_{\chi\ne\chi_0}\abs{\mathbb{E}_{n\sim\mu}\chi(n)}$, taken over the nontrivial Dirichlet
characters modulo $q$, is the largest nontrivial eigenvalue of the weighted Cayley graph
it defines on $(\Z/q\Z)^\times$. Under the generalized Riemann hypothesis (GRH) for Dirichlet
$L$-functions, Jao, Miller and Venkatesan showed that the primes $p\le x$ yield such
expanders with eigenvalues $O(x^{-1/2}\log(qx))$, with unspecified constants.
GRH for Dirichlet $L$-functions now has a machine-checked proof in Lean~4 (Alpöge, with Claude).
Building on that proof, we give fully explicit, unconditional bounds of this type. All of them are
formally verified, with no axioms beyond those of Lean's core.

The main result uses a prime-power measure weighted by a $\cosh$-based test function, and controls it
through the Weil explicit formula and a positivity argument. Its bias is at most an explicit
expression whose leading term is $k_a(0)\log q/(\sqrt{x}\,\int k_a(u)e^{u/2}\,du)$. As a
machine-checked instance: \emph{for every modulus $q\le 2^{64}$ and every nontrivial character $\chi$
modulo $q$, a fixed measure on prime powers $n\le 48400$ has bias at most $1/2$.} This uses a
convolution-square kernel (the $\cosh$ kernel needs $n\le166253$). We also determine the
limit of the method. Among positive-definite kernels the $\cosh$ kernel is optimal, with support-normalized
leading constant $c^*=2/a_*^2\approx1.2236$, where $a_*\approx1.2785$ is the root of $a=1+e^{-a}$. Convolution squares
use the same positivity step and do better. An elementary Fourier identity shows that no kernel gets below $1$, and
this barrier is attained in the limit by the classical weight $\Lambda(n)\log(x/n)$. So the method needs support
$(1+o(1))\log^2q/\varepsilon^2$ for bias $\varepsilon$, and no less.
A cruder Borel--Carath\'eodory bound is also formalized, and we ported a 120-module Lean library
for the Weil explicit formula to the Mathlib version of the GRH proof.
\end{abstract}

\maketitle
\tableofcontents

\section{Introduction}\label{sec:intro}

\subsection{Pseudorandomness from primes}
Fix a modulus $q\ge 1$ and write $\Gq=(\Z/q\Z)^\times$. A probability measure $\mu$ on positive
integers coprime to $q$ pushes forward to a probability measure on $\Gq$. We measure how
pseudorandom it is against the class of \emph{multiplicative characters} of $\Gq$, which are the
nontrivial Dirichlet characters $\chi$ modulo $q$:
\[
  \bias(\mu)=\max_{\chi\ne\chi_0 \bmod q}\Bigl\lvert\sum_n \mu(n)\chi(n)\Bigr\rvert .
\]
This is statistical pseudorandomness for a fixed test class. It is not cryptographic
indistinguishability. It is also a spectral quantity: the characters of $\Gq$ are eigenfunctions of the
averaging operator $Tf(g)=\sum_n\mu(n)f(gn)$ of the weighted Cayley graph $\mathrm{Cay}(\Gq,\mu)$, with
eigenvalues $\sum_n\mu(n)\chi(n)$ (Lemma~\ref{lem:eigen}). Small bias therefore means a uniform
spectral gap. Through the usual consequences of expansion (mixing of random walks, hitting and
Chernoff bounds for expander walks, coset hitting), it yields samplers on $\Gq$ that require
neither the factorization of $q$ nor knowledge of the group structure of $\Gq$.

Natural measures of this kind are supported on primes. Under the generalized Riemann hypothesis for
Dirichlet $L$-functions, the primes $p\le x$ form an expanding generating set of $\Gq$ once $x$ is a
little larger than $(\log q)^2$. This is the case $K=\mathbb{Q}$ of Jao, Miller and Venkatesan
\cite{jmv2009}, whose nontrivial eigenvalues are $O(x^{1/2}\log(xq))$ with absolute but unspecified
constants. Unconditionally, no construction of this kind with polylogarithmic cutoff is known for $\Gq$.
One would imply a polylogarithmic bound for the least character non-residue, which is far beyond
current unconditional methods.

\subsection{The new input}
The Riemann hypothesis and GRH for all Dirichlet $L$-functions have recently been proved in Lean~4
\cite{alpoge2026,grhlean}. The proof has been independently built, replayed by several kernels and
compared against its statement files \cite{sutherlandgrh,sutherlandshortgrh,rubinsteinpage}. The Lean
development is the proof: the accompanying paper reduces the main inequality to stated hypotheses
rather than proving it \cite{rubinsteinpage}. The formal statement we use is
\[
  L(s,\chi)\ne0\quad\text{for all $q\ge1$, all $\chi$ mod $q$, and all } \operatorname{Re}s>\tfrac12,
\]
except at the pole $s=1$ of the principal character. Results that were conditional on GRH for
Dirichlet $L$-functions are thus unconditional, provided one accepts the Lean proof. This covers
neither Hecke $L$-functions of general number fields nor Artin $L$-functions, so the class-group and
isogeny-graph expanders of \cite{jmv2009} remain conditional.

\subsection{Results}
We prove explicit, fully formalized spectral bounds for prime Cayley graphs on $\Gq$.
\begin{enumerate}
\item \textbf{A uniform prime character-sum bound} (Theorem~\ref{thm:A}). For all $q\ge1$, all
  $\chi\ne\chi_0$ and all $N\ge3$,
  \[
    \Bigl\lvert\sum_{0<n<N}\chi(n)\Lambda(n)\bigl(1-\tfrac nN\bigr)\Bigr\rvert
    \le \frac{48e}{\pi}(\log N)^2\sqrt N\,\bigl(\pi\log(2+8q)+J\bigr),
  \]
  where $J=\int_\R\log(4+|t|)(1+t^2)^{-1}\,dt\le6$. The proof uses only nonvanishing of $L(s,\chi)$
  right of the critical line, Borel--Carath\'eodory and one contour shift. It needs no zero counting
  and no explicit formula.
\item \textbf{A sharp explicit bias bound} (Theorems~\ref{thm:sharp} and~\ref{thm:closed}). For a
  $\cosh$-based test function $f_a$ with kernel $k_a=f_a*\tilde f_a$, and weights
  $w(n)=\Lambda(n)k_a(\log n-y)n^{-1/2}$ on units $n\le e^{y+2a}$, we bound the bias by an explicit
  closed-form ratio $\mathrm{num}_B/\mathrm{den}_B$ for all $q$, $a\ge1$ and $y>2a$. The leading term
  is $(a+\sinh a)\log q/(\sqrt{x}\,(\sinh a'+a')^2)$ with $x=e^y$ and $a'=a-1/64$. The proof runs the
  Weil explicit formula for $L(s,\chi)$ and for $\zeta(s)$ through a two-sided positivity argument.
\item \textbf{A machine-checked certificate} (Corollary~\ref{cor:cert}). \emph{For every
  $1\le q\le2^{64}$ and every nontrivial $\chi$ modulo $q$, the measure above with $a=1$ and
  $y=2\log150$ has bias at most $1/2$.} It is supported on prime powers $n\le166253$ coprime to $q$.
  The same guarantee from Theorem~\ref{thm:A} would need $N\approx10^{16}$. With a convolution-square kernel
  (Theorem~\ref{thm:convclosed}, Corollary~\ref{cor:convcert}, also formal) the cutoff drops to $n\le48400$.
\item \textbf{The leading constant of the method} (Propositions~\ref{prop:opt}, \ref{prop:convsq}, \ref{prop:barrier}
  and Corollary~\ref{cor:asymp}; informal). Write the leading term as $c\log q/\sqrt X$, with $X$ the cutoff (top of the support).
  Among positive-definite kernels $k=f*\tilde f$ with $\operatorname{supp}f\subseteq[-a,a]$, the $\cosh$ kernel
  $f\propto\cosh(t/2)\mathbf 1_{[-a,a]}$ is optimal, and optimizing over $a$ gives $c^*=2/a_*^2\approx1.2236$, where
  $a_*\approx1.2785$ is the root of $a=1+e^{-a}$. Convolution squares $f*f$, which are not even, satisfy the same
  positivity step and give $c=1/(1-e^{-L/2})$ on a window of length $L$. Every kernel has $c\ge1$, by the identity
  $\int K(u)e^{u/2}du=\frac1{2\pi}\int\hat K(t)\frac{1/2+it}{1/2-it}dt$ for $K$ supported in $(-\infty,0]$. The value $1$
  is attained in the limit by the classical weight $\Lambda(n)\log(x/n)$. An earlier draft wrongly called $c^*$ optimal
  within the positivity method; see Section~\ref{sec:optimality}.
\end{enumerate}

\subsection{What is new, and what is not}
The shape of Theorem~\ref{thm:A}, and of the leading term in Theorem~\ref{thm:sharp}, is classical
under GRH: compare Bach \cite{bach1990} and Lamzouri--Li--Soundararajan \cite{lls2015} for one-sided
statements, and \cite{jmv2009} for the spectral one. In particular, Lemma~6.1 of \cite{lls2015} already gives the leading term of
Theorem~\ref{thm:sharp} asymptotically (with a $1+o(1)$ factor) for each character, for a general class of kernels. Explicit GRH bounds for
primes in progressions, uniform in $q$, are in \cite{ehp2022}. We claim novelty for the following, with the caveat that a
literature search cannot establish uniqueness:
(i) bias (equivalently, second-eigenvalue) bounds for prime Cayley graphs on $\Gq$ with \emph{fully explicit}
constants and no $o(1)$ terms, uniform in $q$ and in $\chi$ (primitive or not), sharp enough to give small concrete certificates
such as Corollaries~\ref{cor:cert} and~\ref{cor:convcert};
(ii) the exact optimum of the leading constant among positive-definite kernels (Proposition~\ref{prop:opt}), and the
observation that convolution squares reach the limiting constant~$1$ of the method; the barrier itself and the weight attaining it
are classical in substance (Section~\ref{sec:optimality});
(iii) the end-to-end formal verification of (i) and of a concrete certificate, with GRH entering as a
proved theorem rather than an assumption.
The prior results closest to ours are discussed in Section~\ref{sec:related}. In particular, Bach's
constant~$2$ for the least non-residue is already formalized \cite{grhmiller}, and so is the Weil
explicit formula \cite{zeta23}. Our two-sided bounds use the same positivity inequality as
\cite{grhmiller}, in a different (spectral) way.

\subsection{Scope}
We do not improve the exponent $2$ in the $(\log q)^2$ cutoff barrier (Section~\ref{sec:limits}):
doing so is at least as hard as improving the GRH bound for the least prime non-residue. We make no
cryptographic claims. Empirical statistics in Section~\ref{sec:experiments} are evidence about typical
behaviour; they are not proofs.

\section{Setting and notation}\label{sec:background}

\subsection{Bias and Cayley graphs}
A \emph{weighted sample space} is a finite set $S\subset\N$ with weights $w:S\to[0,\infty)$ and mass
$M=\sum_{n\in S}w(n)>0$. For a Dirichlet character $\chi$ modulo $q$ put
\[
  \bias_\chi(S,w)=\frac{\bigl\lvert\sum_{n\in S}w(n)\chi(n)\bigr\rvert}{M}.
\]
If every $n\in S$ is coprime to $q$, then $(S,w)$ defines the probability measure $\mu(g)=M^{-1}\sum_{n\equiv
g}w(n)$ on $\Gq$, and $\bias_\chi(S,w)=\abs{\mathbb E_\mu\chi}$.

\begin{lemma}[eigenvalues; \lean{GRHPseudo.charSum_eq_eigenvalue}]\label{lem:eigen}
Let $T f(g)=\sum_{n\in S}w(n)f(gn)$ for $f:\Z/q\Z\to\C$. For every Dirichlet character $\chi$ modulo
$q$ one has $T\chi=\bigl(\sum_{n\in S}w(n)\chi(n)\bigr)\chi$.
\end{lemma}
\begin{proof}
$T\chi(g)=\sum_n w(n)\chi(g)\chi(n)$, since $\chi$ is completely multiplicative.
\end{proof}

Since the characters of $\Gq$ form an orthogonal basis of functions on $\Gq$, the bias of $\mu$ equals the
second-largest eigenvalue modulus of the normalized weighted Cayley graph $\mathrm{Cay}(\Gq,\mu)$. We
formalize only the eigenvector identity; the completeness statement is standard.

\subsection{Dirichlet $L$-functions and the GRH theorem}
For $\chi$ modulo $q$ we write $L(s,\chi)$ for Mathlib's \lean{DirichletCharacter.LFunction}. It is the
analytic continuation of $\sum_{n\ge1}\chi(n)n^{-s}$ (with $\chi(n)=0$ for $\gcd(n,q)>1$), and it is entire
for $\chi\ne\chi_0$. The GRH theorem we use is the compact edition's \lean{LFunction_ne_zero_of_half_lt_re}
\cite{grhlean}:
\begin{equation}\label{eq:grh}
  \forall q\ge1,\ \forall\chi \bmod q,\ \forall s,\quad \operatorname{Re}s>\tfrac12,\ (\chi\ne\chi_0\text{ or } s\ne1)
  \ \Longrightarrow\ L(s,\chi)\ne0 .
\end{equation}
We also use its consequences \lean{LFunction_eq_zero_re} (zeros with $0<\operatorname{Re}s<1$ lie on the line)
and \lean{riemannHypothesis} (Mathlib's \lean{RiemannHypothesis}).

\subsection{Fourier conventions and the explicit formula}
For $f\in L^1(\R)$ and $z\in\C$ put $\hat f(z)=\int_\R f(u)e^{izu}\,du$ whenever the integral converges, and
$\tilde f(u)=\overline{f(-u)}$. For $f\in C^2_c(\R)$ the Weil test function is $k=f*\tilde f$, i.e.\
$k(u)=\int f(t)\overline{f(t-u)}\,dt$. Then $\hat k(r)=\abs{\hat f(r)}^2\ge0$ for real $r$, and
$k(-u)=\overline{k(u)}$. For a zero $\rho$ we write $\gamma_\rho=(\rho-\tfrac12)/i$, which is real iff
$\operatorname{Re}\rho=\tfrac12$.

We use the explicit formula in the following form, proved in Lean in \cite{zeta23}. For $\chi$ primitive
modulo $q>1$ of parity $\kappa\in\{0,1\}$ and $k\in C^2_c(\R)$, with $m_\rho$ the multiplicity of the
nontrivial zero $\rho$, the zero sum converges absolutely and
\begin{multline}\label{eq:EFchi}
  \sum_\rho m_\rho\,\hat k(\gamma_\rho)\\
  = -\sum_{n\ge1}\frac{\Lambda(n)}{\sqrt n}\Bigl(\chi(n)k(\log n)+\overline{\chi(n)}k(-\log n)\Bigr)\\
  +\frac1{2\pi}\int_\R \hat k(r)\Bigl(\operatorname{Re}\psi\bigl(\tfrac14+\tfrac\kappa2+\tfrac{ir}2\bigr)+\log\tfrac q\pi\Bigr)dr .
\end{multline}
We write $-P_\chi(k)+\Gamma_{q,\kappa}(k)$ for the right side. For $\zeta$,
\begin{multline}\label{eq:EFzeta}
  \sum_\rho m_\rho\,\hat k(\gamma_\rho)=\hat k(\tfrac i2)+\hat k(-\tfrac i2)
  -\sum_{n\ge1}\frac{\Lambda(n)}{\sqrt n}\bigl(k(\log n)+k(-\log n)\bigr)\\
  +\frac1{2\pi}\int_\R\hat k(r)\bigl(\operatorname{Re}\psi(\tfrac14+\tfrac{ir}2)-\log\pi\bigr)dr ,
\end{multline}
which we write as $\mathrm{Pole}(k)-P_\zeta(k)+\Gamma_\zeta(k)$. The Lean statements are
\lean{Zeta23.ThmE.EF_lit_chi_L} and \lean{Zeta23.WeilEF.EF_lit_zetaZeroConfig}; their right sides are the
definitions \lean{literatureRHSchi} and \lean{literatureRHS}.

\subsection{The test function}
Let $S:\R\to[0,1]$ be the smooth transition function, with $S=0$ on $(-\infty,0]$, $S=1$ on $[1,\infty)$
and $S\in C^\infty$ (Mathlib's \lean{Real.smoothTransition}). Following \cite{grhmiller}, put
\[
  \eta_a(t)=S\bigl(64(a-t)\bigr)\,S\bigl(64(a+t)\bigr),\qquad f_a(t)=\cosh(t/2)\,\eta_a(t),\qquad k_a=f_a*\tilde f_a .
\]
Thus $0\le f_a\le\cosh(t/2)$, $f_a(t)=\cosh(t/2)$ for $\abs t\le a-\tfrac1{64}$, and $f_a=0$ for $\abs t\ge a$.
The kernel $k_a$ is real, nonnegative and even, and vanishes for $\abs u>2a$. We write $a'=a-\frac1{64}$ and
\begin{gather*}
  F_\pm(a)=\int_\R f_a(t)e^{\pm t/2}\,dt ,\\
  A(a)=\hat k_a(-\tfrac i2)=F_+(a)F_-(a),\qquad
  B(a)=\hat k_a(\tfrac i2)+\hat k_a(-\tfrac i2)=2F_+(a)F_-(a) .
\end{gather*}

\section{Related work}\label{sec:related}

\paragraph{GRH-based expanders.}
Jao, Miller and Venkatesan \cite{jmv2009} prove under GRH for Hecke $L$-functions that the Cayley graph of a
ray class group with prime-ideal generators of norm at most $x$ has nontrivial eigenvalues
$O(n\,x^{1/2}\log(xq))$, with absolute constants. Over $K=\mathbb Q$ this is the prime Cayley graph on
$\Gq$, and only Dirichlet GRH is needed. Our Theorems~\ref{thm:sharp}--\ref{thm:closed} make this case
fully explicit and formally verified. The class-group and isogeny-graph cases still require GRH for Hecke
$L$-functions of imaginary quadratic fields, which is not covered by the available proof.

\paragraph{Explicit GRH bounds for characters.}
Under GRH, Bach \cite{bach1990} proved that a nontrivial $\chi$ modulo $m$ with $\chi(n)=1$ for all $n<x$
forces $x<2\log^2m$, and Bach and Sorenson \cite{bachsorenson1996} gave explicit bounds for primes in
residue classes. Lamzouri, Li and Soundararajan \cite{lls2015} showed that the least prime outside a proper
subgroup of $\Gq$ is at most $(\log q+B(q))^2$ with $B(q)=o(\log\log q)$. Carneiro, Milinovich,
Quesada-Herrera and Ramos \cite{cmqr2024} optimized the constant for the least quadratic non-residue via
Fourier optimization, obtaining $\limsup n_p/\log^2p<0.7615$. All of these are \emph{one-sided} statements
(existence of a non-residue, or of a prime outside a subgroup). A small bias, by contrast, is a uniform
two-sided statement over all characters at once. The extremal problems differ accordingly
(Section~\ref{sec:optimality}).

Two prior results are closer to ours than these one-sided statements. First, Lemma~6.1 of \cite{lls2015}
(``Asymptotic bounds'') assumes GRH and takes a kernel $K(s)$ with $\abs{K(s)}\ll(1+\abs s)^{-2}$ and
positive inverse Mellin transform $\tilde K$. For non-principal $\chi$ it shows
$\operatorname{Re}\sum_n\Lambda(n)\chi(n)n^{-1/2}\tilde K(n/x)=\theta(1+o(1))\frac{\log q}{2\pi}\int\abs{K(it)}\,dt+O(1+\log q\log x/\sqrt x)$
with $\abs\theta\le1$, and for the principal character it gives $K(\frac12)\sqrt x+O(\cdot)$. In our notation $\tilde K(n/x)=k(\log n-y)$,
$K(\frac12)=\int k(u)e^{u/2}du$ and, for positive-definite $k$, $\frac1{2\pi}\int\abs{K(it)}dt=k(0)$. The leading term of
Theorem~\ref{thm:sharp} is therefore, asymptotically and character by character, that of \cite[Lemma~6.1]{lls2015}.
They use it only for subgroup-escape bounds, with the kernel $K(s)=((e^{\alpha s}-e^{-\alpha s})/s)^2$, i.e.\
$\tilde K(u)=\max(0,2\alpha-\abs{\log u})$; in our notation this is $f=\mathbf 1_{[-\alpha,\alpha]}$, which
Proposition~\ref{prop:opt} shows is not optimal even among positive-definite kernels. Second, Ernvall-Hyt\"onen and Paloj\"arvi \cite{ehp2022} give a fully
explicit GRH bound for $\abs{\psi(x;q,a)-x/\varphi(q)}$, uniform in $q$. Bounding each character this way yields a bias bound
mechanically, but the result is unsmoothed: its leading error is $\asymp\sqrt x\log^2x$ and it has $\sqrt q$-type terms for small $q$, so
the constants are far weaker in the regime $x\approx(\log q)^{2+\delta}$. Neither paper states a bias or eigenvalue bound. Bach's
method, as presented in \cite{mcgowntrevino2021}, weights primes by $\Lambda(n)(n/x)^a\log(x/n)$ with a tuning parameter $a$. Montgomery \cite{montgomery1971} showed under GRH that the least
non-residue is $\Omega(\log q\log\log q)$ infinitely often; Graham and Ringrose \cite{grahamringrose1990}
proved $\Omega(\log q\log\log\log q)$ unconditionally.

\paragraph{Formalizations built on the GRH proof.}
Bach's theorem with constant $2$ and Miller's primality test with Bach's bound are formalized in
\cite{grhmiller}, using the Weil explicit formula of \cite{zeta23} and the positivity inequality that we
also use. Sutherland's verification of Schoenfeld's RH bound \cite{schoenfeldlean}, and the
rh-prime-error project \cite{rhprimeerror}, concern $\zeta$. We reuse code from all of these
(Section~\ref{sec:formal}).

\paragraph{Small-bias spaces and Cayley expanders.}
Random multisets of $O(\log\abs G/\varepsilon^2)$ elements are $\varepsilon$-biased in any finite group
\cite{alonroichman1994}. Explicit near-optimal constructions exist for $\mathbb F_2^n$ \cite{tashma2017}. For
explicitly given abelian groups, Jalan and Moshkovitz \cite{jalanmoshkovitz2021} construct near-optimal
expanders deterministically in time polynomial in $\abs G$. These constructions need the group structure
and run in time $\operatorname{poly}(\abs G)=\operatorname{poly}(q)$. Prime-based sets need neither the
factorization of $q$ nor a description of $\Gq$, and their elements are small integers independent of
$q$; the price is support $\asymp(\log q)^2/\varepsilon^2$ rather than $\log q/\varepsilon^2$. For
background on expanders see \cite{hlw2006}; for small-bias spaces, \cite{naornaor1993}.

\paragraph{Algorithmic consequences already in the literature.}
Squareness of integers given by straight-line programs is in coRP under GRH \cite{gaillardjindal2023}.
As far as we can see, the GRH needed there is only for quadratic Dirichlet characters, so this becomes
unconditional; we have not checked every detail.
Identity testing for radical expressions with square roots \cite{bnsw2022} uses GRH for multiquadratic
fields, whose Dedekind zeta functions are products of Dirichlet $L$-functions. We do not claim these;
they illustrate that the relevant input is exactly~\eqref{eq:grh}.

\section{Main results}\label{sec:results}

Each statement below is labelled with the Lean declaration that proves it. ``Formal'' means machine-checked,
with axioms \lean{propext}, \lean{Classical.choice} and \lean{Quot.sound} only, and with
\eqref{eq:grh} discharged by the GRH proof (Section~\ref{sec:formal}). Statements marked ``informal'' are
proved in this paper but are not formalized.

\subsection{A uniform bound via Borel--Carath\'eodory}

\begin{theorem}[formal; \lean{GRHPseudo.Main.primeCharSum}]\label{thm:A}
For every $q\ge1$, every $\chi\ne\chi_0$ modulo $q$ (primitive or not) and every integer $N\ge3$,
\[
  \Bigl\lvert\sum_{0<n<N}\chi(n)\Lambda(n)\Bigl(1-\frac nN\Bigr)\Bigr\rvert\le\frac{48e}{\pi}(\log N)^2\sqrt N\,
  \bigl(\pi\log(2+8q)+J\bigr),\qquad J=\int_\R\frac{\log(4+\abs t)}{1+t^2}\,dt .
\]
\end{theorem}

\begin{corollary}[formal; \lean{GRHPseudo.Main.bias_primeMeasure}]\label{cor:A}
Let $S$ be the set of $n\in(0,N)$ coprime to $q$, with weights $\Lambda(n)(1-n/N)$. For $N\ge3$ with
$L:=\frac12\bigl(\lfloor N/2\rfloor\log2-\log(\lfloor N/2\rfloor+1)\bigr)-\log q\,\lfloor\log_2N\rfloor>0$, the
bias of $S$ against every $\chi\ne\chi_0$ is at most the right side of Theorem~\ref{thm:A} divided by $L$.
\end{corollary}

The constants here are poor: to certify bias $\le1/2$ for $q\le2^{64}$ one needs $N\approx9.7\cdot10^{15}$
(Table~\ref{tab:params}). The loss is intrinsic to the route. Borel--Carath\'eodory bounds $L'/L$
pointwise on $\operatorname{Re}s=\frac12+1/\log N$ only by $O((\log N)^2\log q)$, whereas the explicit formula
avoids pointwise bounds altogether.

\subsection{The sharp bound}

Fix $a\ge0$ and $y>2a$, put $x=e^y$ and $M=\lfloor e^{y+2a}\rfloor+1$, and define the sample space
\[
  S_q(a,y)=\{0\le n<M:\gcd(n,q)=1\},\qquad w_{a,y}(n)=\frac{\Lambda(n)\,k_a(\log n-y)}{\sqrt n}.
\]
Its weights are nonnegative and supported on prime powers in $[xe^{-2a},xe^{2a}]$. With $f=f_a$, put
\begin{align*}
  I_a(\beta)&=\int_\R\abs{\hat f_a(r)}^2\bigl(\log(4+\abs r)+\beta\bigr)\,dr,\qquad
  \mathrm{SP}(a)=\sum_{n\le e^{2a}}\frac{\Lambda(n)}{\sqrt n}\bigl(\abs{k_a(\log n)}+\abs{k_a(-\log n)}\bigr),\\
  U_q(a,y)&=\sum_{n<M,\ \gcd(n,q)>1}\abs{w_{a,y}(n)} ,
\end{align*}
and
\begin{align*}
  \mathrm{Num}&=k_a(0)\log\tfrac q\pi+\mathrm{SP}(a)+\frac{I_a(5)+I_a(\frac12)}{2\pi}+U_q(a,y),\\
  \mathrm{Den}&=e^{y/2}A(a)-B(a)-\frac{I_a(5+\log\pi)+I_a(\frac12)}{2\pi}-U_q(a,y).
\end{align*}

\begin{theorem}[formal; \lean{GRHPseudo.Sharp.bias_sharp}, \lean{sharp_charSum_mass}]\label{thm:sharp}
For every $q\ge1$, every $\chi\ne\chi_0$ modulo $q$, every $a\ge0$ and every $y>2a$:
\[
  \Bigl\lvert\sum_{n\in S_q(a,y)}w_{a,y}(n)\chi(n)\Bigr\rvert\le\mathrm{Num},\qquad
  \sum_{n\in S_q(a,y)}w_{a,y}(n)\ge\mathrm{Den}.
\]
In particular, if $\mathrm{Den}>0$ then $\bias_\chi(S_q(a,y),w_{a,y})\le\mathrm{Num}/\mathrm{Den}$.
\end{theorem}

The leading term of $\mathrm{Num}/\mathrm{Den}$ is $k_a(0)\log q/(\sqrt x\,A(a))$. All remaining terms are
either independent of $q$ or $O(\log q\log x/\sqrt x)$.

\begin{remark}
This leading term agrees with the asymptotic Lemma~6.1 of Lamzouri--Li--Soundararajan \cite{lls2015} (see
Section~\ref{sec:related}). Theorem~\ref{thm:sharp} adds the following:
\begin{itemize}
\item every lower-order term is explicit, with no $o(1)$;
\item the numerator bounds $\abs{\sum w\chi}$ itself, not just its real part, and holds for all $\chi\ne\chi_0$, primitive or not;
\item the mass bound is explicit, obtained from the explicit formula for $\zeta$ under RH;
\item the whole statement is machine-checked.
\end{itemize}
The explicit form makes small certificates possible: with an $o(1)$ one cannot conclude anything at any specific $q$.
\end{remark}

\subsection{Closed form}
Let $K(a)=\sinh a+a\cosh a$, $C(a)=(4\sinh\frac a2)^2+(2\sinh\frac a2+2\cosh\frac a2)^2$,
$U_B=K(a)e^{a-y/2}\log q\,\lfloor\log_2M\rfloor$, and
\begin{align*}
  \mathrm{num}_B&=(\sinh a+a)\log q+\frac{4K(a)}{e}\bigl(\lfloor e^{2a}\rfloor+1\bigr)
     +\frac{C(a)\bigl(12+\frac{11}2\pi\bigr)}{2\pi}+U_B,\\
  \mathrm{den}_B&=e^{y/2}\bigl(\sinh a'+a'\bigr)^2-2(\sinh a+a)^2-\frac{C(a)\bigl(12+(\frac{11}2+\log\pi)\pi\bigr)}{2\pi}-U_B .
\end{align*}

\begin{theorem}[formal; \lean{GRHPseudo.Main.bias_sharp_closed}]\label{thm:closed}
For every $q\ge1$, $\chi\ne\chi_0$ modulo $q$, $a\ge1$ and $y>2a$ with $\mathrm{den}_B>0$,
\[
  \bias_\chi\bigl(S_q(a,y),w_{a,y}\bigr)\le\frac{\mathrm{num}_B}{\mathrm{den}_B}.
\]
\end{theorem}

\begin{corollary}[formal; \lean{GRHPseudo.Main.bias_le_half_of_le_two_pow_64}]\label{cor:cert}
Let $a=1$ and $y=2\log150$, so $x=22500$ and $S_q(1,y)$ consists of the integers below
$\lfloor22500e^2\rfloor+1=166254$ coprime to $q$. For every $1\le q\le2^{64}$ and every nontrivial
Dirichlet character $\chi$ modulo $q$,
\[
  \bias_\chi\bigl(S_q(1,y),w_{1,y}\bigr)\le\tfrac12 .
\]
Equivalently, every nontrivial eigenvalue of the weighted Cayley graph $\mathrm{Cay}(\Gq,w_{1,y})$ has modulus
at most $1/2$. (The proof gives $\mathrm{num}_B\le236.6$ and $\mathrm{den}_B\ge555.8$, hence bias $\le0.4257$.)
\end{corollary}

\subsection{The leading constant}

\begin{proposition}[informal; abstract core formal: \lean{GRHPseudo.inner_mul_inner_le}]\label{prop:opt}
Let $a>0$ and let $f\in L^2[-a,a]$ be real and nonzero, with $k=f*\tilde f$. Then
\[
  \frac{k(0)}{\int_\R k(u)e^{u/2}\,du}\ \ge\ \frac1{a+\sinh a},
\]
with equality if and only if $f$ is a nonzero multiple of $\cosh(t/2)\mathbf 1_{[-a,a]}$. By the
Boas--Kac--Krein factorization \cite{boaskac1945,ehmgneitingrichards2004}, every continuous positive-definite
$k$ supported in $[-2a,2a]$ is of this form.
\end{proposition}

The quantity in Proposition~\ref{prop:opt} is the coefficient of $\log q/\sqrt x$ that the positivity method
produces from a \emph{positive-definite} kernel supported in $[-2a,2a]$ (Section~\ref{sec:optimality}). It is not
the limit of the method: non-even kernels whose Fourier transform is dominated by a positive-definite one do better.
Normalizing by the top $X$ of the support:

\begin{corollary}[informal]\label{cor:asymp}
Write the leading term of the positivity bound, or of the bound of \cite[Lemma~6.1]{lls2015}, as $c\log q/\sqrt X$, where
$X$ is the top of the support.
\begin{enumerate}
\item \emph{(Barrier.)} Every kernel gives $c\ge1$ (Proposition~\ref{prop:barrier}).
\item \emph{(Positive-definite kernels.)} A positive-definite kernel supported in an interval of length $4a$ gives
  $c\ge e^a/(a+\sinh a)\ge c^*=2/a_*^2=1.22363\ldots$, where $a_*=1.27846\ldots$ is the root of $a=1+e^{-a}$.
\item \emph{(Achievability.)} For every $\delta>0$ there is a smooth nonnegative convolution-square kernel with
  $\bias\le(1+\delta+o(1))\log q/\sqrt X$ as $q\to\infty$ and $X/\log^2q\to\infty$. Hence support
  $X\ge(1+o(1))\log^2q/\varepsilon^2$ suffices for bias $\le\varepsilon$, and within the method nothing smaller does.
\end{enumerate}
For the formally verified kernel $f_a$ (positive-definite, transition width $1/64$), Theorem~\ref{thm:closed} gives leading
coefficient at most $e^a(\sinh a+a)/(\sinh a'+a')^2$, which equals $1.2626$ at $a=a_*$ and $1.2964$ at $a=1$.
\end{corollary}

\subsection{An explicit bound with a convolution-square kernel}
Corollary~\ref{cor:asymp} says that a non-even kernel beats every positive-definite one. We have made this explicit
and formal for one family. Let $A>0$, let $S$ be Mathlib's \lean{Real.smoothTransition}, and put
\[
  f_A(t)=e^{t/2}\,S\bigl(64(t+A)\bigr)\,S(-64t),\qquad K_A=f_A*f_A,\qquad M_A=f_A*\tilde f_A .
\]
So $f_A\ge0$ is supported in $[-A,0]$ and equals $e^{t/2}$ on $[-A+\frac1{64},-\frac1{64}]$, $K_A\ge0$ is supported in
$[-2A,0]$, and $\abs{\hat K_A}=\abs{\hat f_A}^2=\hat M_A$ on $\R$ (Proposition~\ref{prop:convsq}). For $y>2A$ and $x=e^y$
let $S^\square_q(y)=\{0\le n\le\lfloor x\rfloor:\gcd(n,q)=1\}$ with weights
$w^\square_{A,y}(n)=\Lambda(n)K_A(\log n-y)n^{-1/2}$, which are supported on prime powers in $[xe^{-2A},x]$. For $R>0$
put $G(B)=(\log(4+R)+B)(1-e^{-A})+36/(\pi R)$ and $U^\square_B=Ae^{-y/2}\log q\,\lfloor\log_2(\lfloor x\rfloor+1)\rfloor$, and
\begin{align*}
  \mathrm{num}^\square_B&=(1-e^{-A})\log q+\tfrac2e\bigl(\lfloor e^{A}\rfloor+1\bigr)+G(5)+G(\tfrac12)+U^\square_B,\\
  \mathrm{den}^\square_B&=e^{y/2}\bigl(e^{-1/64}-e^{-A+1/64}\bigr)^2-2A-G(5+\log\pi)-G(\tfrac12)-U^\square_B .
\end{align*}

\begin{theorem}[formal; \lean{GRHPseudo.Main.bias_convSq_closed}]\label{thm:convclosed}
For every $q\ge1$, $\chi\ne\chi_0$ modulo $q$, $A\ge\frac1{32}$, $R>0$ and $y>2A$ with $\mathrm{den}^\square_B>0$,
\[
  \bias_\chi\bigl(S^\square_q(y),w^\square_{A,y}\bigr)\le\frac{\mathrm{num}^\square_B}{\mathrm{den}^\square_B}.
\]
\end{theorem}

The proof runs the argument of Theorem~\ref{thm:sharp} with the zero sum of the twisted $K_A$ bounded by that of $M_A$
(\lean{norm_primeSide_twist_le_maj}, \lean{primeSideZeta_twist_eq_maj}); the intermediate bound with exact constants is
\lean{GRHPseudo.Main.bias_convSq_sharp}. It also uses a sharper gamma-side estimate (\lean{Ilog_div_le_sharp}): if
$f\in C^2_c$ and $\abs r\,\abs{\hat f(r)}\le C$, then for $B\ge0$ and $R>0$
\[
  \frac1{2\pi}\int_\R\abs{\hat f(r)}^2\bigl(\log(4+\abs r)+B\bigr)dr\le\bigl(\log(4+R)+B\bigr)(f*\tilde f)(0)+\frac{4C^2}{\pi R}.
\]
Here $C=3$, from $\int\abs{f_A'}\le3$. The leading coefficient of $\mathrm{num}^\square_B/\mathrm{den}^\square_B$, normalized by the top
$x$ of the support, is $(1-e^{-A})/(e^{-1/64}-e^{-A+1/64})^2$, which equals $1.162$ at $A=\frac94$; for the formal $\cosh$ kernel
the corresponding value at $a=1$ is $1.2964$.

\begin{corollary}[formal; \lean{GRHPseudo.Main.bias_le_half_of_le_two_pow_64_convSq}]\label{cor:convcert}
Let $A=\frac94$, $R=12$ and $y=2\log220$, so $x=48400$ and $S^\square_q(y)$ consists of the integers $n\le48400$ coprime to $q$.
For every $1\le q\le2^{64}$ and every nontrivial Dirichlet character $\chi$ modulo $q$,
\[
  \bias_\chi\bigl(S^\square_q(y),w^\square_{9/4,y}\bigr)\le\tfrac12 .
\]
(The proof gives $\mathrm{num}^\square_B\le65.7$ and $\mathrm{den}^\square_B\ge143.9$, hence bias $\le0.4566$.) The support is $3.4$ times
smaller than in Corollary~\ref{cor:cert}.
\end{corollary}

\section{Proofs}\label{sec:proofs}

Each proof below follows the Lean development. We name the modules (under
\lean{lean/GRHPseudo/AnalyticNumberTheory/}) so the reader can compare.

\subsection{Theorem~\ref{thm:A}: Borel--Carath\'eodory and one contour shift}
\emph{Step 1: growth of $L$.} For $\chi\ne\chi_0$ modulo $N$ and $\operatorname{Re}s\ge\frac12$,
$\abs{L(s,\chi)}\le1+4N\abs s$ (\lean{norm_LFunction_le_linear}). The proof is partial summation: the partial sums of
$\chi$ are bounded by $2N$. The resulting integral is analytic for $\operatorname{Re}s>0$, and the identity
theorem extends the formula from $\operatorname{Re}s>1$. This is adapted from the readable GRH edition.

\emph{Step 2: the GRH step.} Since $L(c,\chi)L(c,\chi\mu)=1$ and $\abs{L(c,\chi\mu)}\le\zeta(2)\le2$, we
have $\abs{L(c,\chi)}\ge\frac12$ at $c=2+it$. On the disc about $c$ of radius $\frac32-\frac\varepsilon2$,
$L$ has no zeros by \eqref{eq:grh}. The Borel--Carath\'eodory bound for logarithmic derivatives then gives, for
$0<\varepsilon\le1$ and $\frac12+\varepsilon\le\operatorname{Re}s\le2$,
\[
  \Bigl\lvert\frac{L'}{L}(s,\chi)\Bigr\rvert\le\frac{48}{\varepsilon^2}\log\bigl(2(1+4q(\abs{\operatorname{Im}s}+4))\bigr)
  \qquad(\text{\lean{norm_logDeriv_LFunction_le}}).
\]

\emph{Step 3: Perron and shift.} The weighted Perron formula with kernel $N^s(1/s-1/(s+1))$ (from
\cite{schoenfeldlean}) expresses the sum in Theorem~\ref{thm:A} as a vertical integral of
$-L'/L(s)\,N^s(1/s-1/(s+1))$ on $\operatorname{Re}s=2$. We shift it to $\operatorname{Re}s=\frac12+1/\log N$
with a rectangle-contour lemma (\lean{integral_vertical_eq}); the horizontal edges decay like $1/\abs T$.
Then we bound the integrand using $\abs{1/s-1/(s+1)}\le2/(1+t^2)$ and $N^{1/2+1/\log N}=e\sqrt N$.

\emph{Corollary~\ref{cor:A}.} The mass is at least $\psi(\lfloor N/2\rfloor)/2$ (Mathlib's \lean{Chebyshev.psi_ge}).
Removing non-units costs at most $\log q\,\lfloor\log_2N\rfloor$, since each non-unit prime power $p^k<N$ has
$p\mid q$ and $k\le\log_2N$.

\subsection{Theorem~\ref{thm:sharp}}
\emph{Two-sided positivity} (\lean{Sharp/TwoSided.lean}).
\begin{lemma}\label{lem:twosided}
Let $\chi$ be primitive modulo $q>1$ of parity $\kappa$, let $f\in C^2_c(\R)$, $k=f*\tilde f$, $y\in\R$, and
$k_y=k(\cdot-y)$. Then
\[
  \abs{P_\chi(k_y)}\le\abs{\Gamma_{q,\kappa}(k_y)}+\operatorname{Re}\bigl(\Gamma_{q,\kappa}(k)-P_\chi(k)\bigr).
\]
\end{lemma}
\begin{proof}
We have $\hat k_y(z)=e^{izy}\hat k(z)$. By GRH every $\gamma_\rho$ is real, so each term of the zero sum
in \eqref{eq:EFchi} for $k_y$ equals $m_\rho e^{i\gamma_\rho y}\abs{\hat f(\gamma_\rho)}^2$. The zero sum
for $k_y$ therefore has modulus at most $S=\sum_\rho m_\rho\abs{\hat f(\gamma_\rho)}^2$, which is the
zero sum for $k$ itself. Apply \eqref{eq:EFchi} to both $k_y$ and $k$.
\end{proof}

\emph{Finite sums} (\lean{Sharp/FiniteSums.lean}). Since $k_a(u)=0$ for $\abs u>2a$ and $y>2a$, the conjugate
leg $k_a(-\log n-y)$ vanishes for every $n$. The prime side of $k_{a,y}$ is then the finite sum
$\sum_{n<M}\Lambda(n)n^{-1/2}\chi(n)k_a(\log n-y)$.

\emph{Gamma sides} (\lean{Sharp/GammaBounds.lean}). We need the following facts.
\begin{enumerate}
\item $\abs{\operatorname{Re}\psi(\frac14+\frac\kappa2+\frac{ir}2)}\le\log(4+\abs r)+5$. The upper side is from
  \cite{grhmiller}. The lower side uses $\operatorname{Re}\psi(\alpha+it)=-\gamma-\alpha/(\alpha^2+t^2)+\sum_n(\ldots)$ with
  nonnegative terms, which is proved in \cite{zeta23}, and $\gamma<\frac23$.
\item \emph{The $\log q$ term vanishes for the twisted kernel}: by Fourier inversion at $-y$,
  $\frac1{2\pi}\int\hat k_y(r)\,dr=k(-y)=0$. Hence $\abs{\Gamma_{q,\kappa}(k_y)}\le I(5)/(2\pi)$. This is what keeps the
  coefficient of $\log q$ equal to $k(0)$ rather than $2k(0)$.
\item Plancherel, $\frac1{2\pi}\int\abs{\hat f}^2=k(0)$, gives
  $\operatorname{Re}\Gamma_{q,\kappa}(k)\le k(0)\log(q/\pi)+I(\frac12)/(2\pi)$.
\item For $\zeta$: $\abs{\Gamma_\zeta(k_y)}\le I(5+\log\pi)/(2\pi)$ and $\operatorname{Re}\Gamma_\zeta(k)\le I(\frac12)/(2\pi)$.
\end{enumerate}

\emph{Main term} (\lean{Sharp/MainTerm.lean}). By the same argument with \eqref{eq:EFzeta} and RH,
$P_\zeta(k_y)=\mathrm{Pole}(k_y)+\Gamma_\zeta(k_y)-E$ with $\abs E\le\operatorname{Re}(\mathrm{Pole}(k)-P_\zeta(k)+\Gamma_\zeta(k))$.
Moreover $\mathrm{Pole}(k_y)=e^{-y/2}\hat k(\frac i2)+e^{y/2}\hat k(-\frac i2)$. Since $k_a\ge0$, both
$\operatorname{Re}\hat k_a(\frac i2)\ge0$ and $\operatorname{Re}P_\zeta(k_a)\ge0$.

\emph{Assembly} (\lean{Sharp/Assembly.lean}). Let $\chi^*$ be the primitive character modulo the conductor
$q^*\mid q$ that induces $\chi$; then $q^*>1$ since $\chi\ne\chi_0$. On units modulo $q$ we have $\chi=\chi^*$,
so the character sum over $S_q(a,y)$ differs from $P_{\chi^*}(k_{a,y})$ by at most $U_q(a,y)$. We also use
$\log(q^*/\pi)\le\log(q/\pi)$ and $-\operatorname{Re}P_{\chi^*}(k_a)\le\mathrm{SP}(a)$. The mass over all $n<M$ is
$\operatorname{Re}P_\zeta(k_{a,y})$, and the mass over units is smaller by at most $U_q$. Combining gives the two
inequalities of Theorem~\ref{thm:sharp}.

\subsection{Theorem~\ref{thm:closed} and Corollary~\ref{cor:cert}}
We need closed-form bounds on the constants of Theorem~\ref{thm:sharp} (\lean{Sharp/KernelBounds.lean},
\lean{Sharp/ConstBounds.lean}).
\begin{enumerate}
\item $\abs{k_a(u)}\le\sinh(a-\frac u2)+(a-\frac u2)\cosh\frac u2\le K(a)$ for $0\le u\le2a$ (from \cite{grhmiller}); by symmetry
  this holds for all $\abs u\le2a$. Also $k_a(0)\le\sinh a+a$.
\item $A(a)=F_+F_-$ and $B(a)=2F_+F_-$, from $\widehat{f*\tilde f}(z)=\hat f(z)\overline{\hat f(\bar z)}$ at $z=\mp\frac i2$.
  Since $f_a=\cosh(t/2)$ on $[-a',a']$ and $0\le f_a\le\cosh(t/2)\mathbf 1_{[-a,a]}$, and
  $\int_{-c}^c\cosh(t/2)e^{\pm t/2}dt=\sinh c+c$, we get $\sinh a'+a'\le F_\pm\le\sinh a+a$.
\item $I_a(\beta)\le C(a)(6+\beta\pi)$, from $\abs{\hat f_a}\le4\sinh\frac a2$, $\abs{r\hat f_a(r)}\le2\sinh\frac a2+2\cosh\frac a2$
  and $J\le6$ (all from \cite{grhmiller}).
\item $\Lambda(n)/\sqrt n\le\log n/\sqrt n\le2/e$, which gives $\mathrm{SP}(a)\le4K(a)e^{-1}(\lfloor e^{2a}\rfloor+1)$.
\item On the support, $\sqrt n\ge e^{(y-2a)/2}$. Together with the non-unit prime-power bound this gives $U_q\le U_B$.
\end{enumerate}
For Corollary~\ref{cor:cert} (\lean{Sharp/Certificate64.lean}) take $a=1$ and $y=2\log150$, so $K(1)=e$ and
$e^{y/2}=150$. Mathlib's certified bounds $2.7182818283<e<2.7182818286$, $3.141592<\pi<3.141593$ and
$\log2<0.6931471808$ then give the following.
\begin{itemize}
\item $\log q\le44.3614197$.
\item $e^{1/2}<1.6488$, hence $C(1)\le15.219$.
\item $e^{63/64}\ge e\cdot\frac{63}{64}$ and $e^{-63/64}\le\frac{64}{63}e^{-1}$, hence $(\sinh a'+a')^2\ge4.5599$.
\item $\lfloor e^2\rfloor\le7$, and $\lfloor22500e^2\rfloor+1<2^{18}$, so $\lfloor\log_2M\rfloor\le17$ and $U_B\le37.15$.
\end{itemize}
Altogether $\mathrm{num}_B\le236.6$ and $\mathrm{den}_B\ge555.8$.

\section{The leading constant of the method}\label{sec:optimality}

\emph{Correction (2026-11-27).} An earlier draft of this paper stated that the $\cosh$ kernel is optimal
``within the positivity method'' and that no kernel gives a leading constant below $c^*\approx1.2236$ there.
That was wrong as framed. Proposition~\ref{prop:opt} is correct, but only for positive-definite kernels.
Convolution squares (Proposition~\ref{prop:convsq}) use the same positivity step and do better. The sharp
limit of the method is the constant~$1$ (Proposition~\ref{prop:barrier}).

\subsection{Where the leading coefficient comes from}
Let $K$ be a real kernel and suppose $M$ is a positive-definite function with $\abs{\hat K}\le\hat M$ on $\R$.
Under GRH the zeros $\frac12+i\gamma$ have $\gamma$ real, so the twist $K_y=K(\cdot-y)$ satisfies
$\abs{\sum_\gamma\hat K_y(\gamma)}\le\sum_\gamma\hat M(\gamma)$, and the right side is evaluated by the explicit
formula for $M$. This is the positivity step. For $K=f*\tilde f$ we may take $M=K$; this is the case formalized in
Theorem~\ref{thm:sharp}. In general, only two terms of the resulting bound are large. In the numerator the only
term growing with $q$ is $M(0)\log(q/\pi)$, from the gamma factor through Plancherel. In the denominator the only
term growing with $x$ is $\sqrt x\int K(u)e^{u/2}du$. As $q\to\infty$ and $x/\log^2q\to\infty$ with $K$ fixed,
the method therefore gives
\begin{equation}\label{eq:leading}
  \bias\le\bigl(\rho+o(1)\bigr)\frac{\log q}{\sqrt x},\qquad
  \rho=\frac{M(0)}{\int K(u)e^{u/2}\,du}\ \ge\ \frac{\frac1{2\pi}\int\abs{\hat K}}{\int K(u)e^{u/2}\,du}.
\end{equation}
The last expression is also the leading coefficient in \cite[Lemma~6.1]{lls2015}, which bounds the zero sum by
$\sum_\gamma\abs{\hat K(\gamma)}$ and zero counting instead of by positivity. For $K=f*\tilde f$ with $f$ real,
$\rho=\norm f_2^2/(\langle f,e^{t/2}\rangle\langle f,e^{-t/2}\rangle)$, since $K(0)=\norm f^2$ and
$\hat K(z)=\hat f(z)\overline{\hat f(\bar z)}$. If $\operatorname{supp}K\subseteq(-\infty,T]$, the measure is
supported on $n\le X=xe^{T}$, and the support-normalized constant is $c=e^{T/2}\rho$: then
$\bias\le(c+o(1))\log q/\sqrt X$.

\subsection{Positive-definite kernels}
\begin{lemma}[formal; \lean{GRHPseudo.inner_mul_inner_le}]\label{lem:ranktwo}
In a real inner product space, for all $f,\alpha,\beta$,
\[
  \langle f,\alpha\rangle\langle f,\beta\rangle\le\tfrac12\bigl(\norm\alpha\norm\beta+\langle\alpha,\beta\rangle\bigr)\norm f^2 .
\]
If $\norm\alpha\norm\beta>0$, equality holds iff $f$ is a multiple of $\alpha/\norm\alpha+\beta/\norm\beta$.
\end{lemma}
\begin{proof}
Put $u=\alpha/\norm\alpha$ and $v=\beta/\norm\beta$. Then
$4\langle f,u\rangle\langle f,v\rangle=\langle f,u+v\rangle^2-\langle f,u-v\rangle^2\le\norm f^2\norm{u+v}^2=2(1+\langle u,v\rangle)\norm f^2$.
Equality requires $f\perp u-v$ and $f\parallel u+v$; note that $u+v\perp u-v$.
\end{proof}

\begin{proof}[Proof of Proposition~\ref{prop:opt}]
Work in $L^2[-a,a]$ with $\alpha=e^{t/2}$ and $\beta=e^{-t/2}$. Then $\norm\alpha^2=\norm\beta^2=\int_{-a}^ae^t\,dt=2\sinh a$ and
$\langle\alpha,\beta\rangle=2a$. Lemma~\ref{lem:ranktwo} gives $\rho\ge1/(\sinh a+a)$, with equality iff
$f\propto e^{t/2}+e^{-t/2}=2\cosh(t/2)$ on $[-a,a]$. The factorization statement is the theorem of Boas and Kac
and of Krein: a continuous positive-definite function supported in $[-2a,2a]$ is $f*\tilde f$ for some
$f\in L^2[-a,a]$ \cite{boaskac1945,ehmgneitingrichards2004}. Hence the bound covers every positive-definite
kernel, i.e.\ every kernel that is its own majorant.
\end{proof}

A positive-definite kernel supported in $[-2a,2a]$ has $T=2a$, so its support-normalized constant is at least
$e^a/(a+\sinh a)$, whose minimum over $a$ is $c^*=2/a_*^2\approx1.2236$ (Corollary~\ref{cor:asymp}). Positive-definite
kernels are even. Since the main term weights the window by $e^{u/2}$, an even kernel wastes the lower half of it.

\subsection{Convolution squares}
\begin{proposition}[informal]\label{prop:convsq}
Let $f\ge0$ be real, integrable and supported in $[T/2-L/2,\,T/2]$, and let $K=f*f$. Then $K\ge0$, $\operatorname{supp}K\subseteq[T-L,T]$,
and $M=f*\tilde f$ is a positive-definite majorant: $\abs{\hat K}=\abs{\hat f}^2=\hat M$. Hence the positivity step applies to
$K$ unchanged, with
\[
  \rho=\frac{\norm f_2^2}{\langle f,e^{t/2}\rangle^2},\qquad
  c=e^{T/2}\rho\ \ge\ \frac1{1-e^{-L/2}},
\]
with equality iff $f\propto e^{t/2}$ on its support interval.
\end{proposition}
\begin{proof}
$\hat K=\hat f^2$ and $\hat M=\abs{\hat f}^2$. Next, $M(0)=\norm f^2$ and $\int K(u)e^{u/2}du=\hat K(\frac i2)=(\int f(t)e^{t/2}dt)^2$.
By Cauchy--Schwarz, $\langle f,e^{t/2}\rangle^2\le\norm f^2\int_{T/2-L/2}^{T/2}e^t\,dt=\norm f^2e^{T/2}(1-e^{-L/2})$, with equality iff
$f\propto e^{t/2}$ there.
\end{proof}

At window length $L=4a$, the length of a positive-definite kernel supported in $[-2a,2a]$, this improves on the $\cosh$ kernel:
$1/(1-e^{-2a})$ against $e^a/(a+\sinh a)$, i.e.\ $1.157$ against $1.250$ at $a=1$, $1.084$ against $1.224$ at $a=a_*$, and $1.019$
against $1.313$ at $a=2$. As $L\to\infty$ the constant tends to~$1$, while the positive-definite constant grows. In the
proof of Theorem~\ref{thm:sharp} only the two-sided zero-sum estimate changes: the zero sum of $K_y$ is bounded by that
of $M$ rather than by its own. Theorem~\ref{thm:convclosed} and Corollary~\ref{cor:convcert} are the formal, explicit
version for a smoothed $f=e^{t/2}\mathbf 1_{[-A,0]}$; the Cauchy--Schwarz optimality statement itself is not formalized.

\subsection{The barrier}
\begin{proposition}[informal]\label{prop:barrier}
Let $K\in L^1(\R)$ be real and continuous, with $\hat K\in L^1(\R)$ and $\operatorname{supp}K\subseteq(-\infty,0]$. Then
\[
  \int_\R K(u)e^{u/2}\,du=\frac1{2\pi}\int_\R\hat K(t)\,\frac{\frac12-it}{\frac12+it}\,dt,
  \qquad\text{so}\qquad
  \Bigl\lvert\int_\R K(u)e^{u/2}\,du\Bigr\rvert\le\frac1{2\pi}\int_\R\abs{\hat K(t)}\,dt .
\]
Equality holds for $K(u)=\abs u\,e^{u/2}\mathbf 1_{u\le0}$. Consequently every kernel and every majorant in
\eqref{eq:leading} give support-normalized constant $c\ge1$.
\end{proposition}
\begin{proof}
By Fourier inversion, $K(u)=\frac1{2\pi}\int\hat K(t)e^{itu}dt$ for all $u$. Integrating against $e^{u/2}$ over $u\le0$ and
using Fubini (justified by $\int_{-\infty}^0e^{u/2}du<\infty$ and $\hat K\in L^1$) gives
$\int_{-\infty}^0K(u)e^{u/2}du=\frac1{2\pi}\int\hat K(t)(\frac12+it)^{-1}dt$. Since $K$ is continuous and vanishes on
$(0,\infty)$, $0=K(0)=\frac1{2\pi}\int\hat K(t)\,dt$. Subtracting, and using
$(\frac12+it)^{-1}-1=(\frac12-it)/(\frac12+it)$, which has modulus~$1$, proves the identity and the inequality. For
$K=\abs ue^{u/2}\mathbf 1_{u\le0}$ we have $\hat K(t)=(\frac12-it)^{-2}$, so the integrand is $(\frac14+t^2)^{-1}>0$, and both sides
equal~$1$. For the consequence, translate a kernel supported in $(-\infty,T]$ to $(-\infty,0]$; then
$\int K(u)e^{u/2}du\le e^{T/2}\frac1{2\pi}\int\abs{\hat K}\le e^{T/2}M(0)$, i.e.\ $c\ge1$.
\end{proof}

The extremal kernel is itself a convolution square, $e*e$ with $e(t)=e^{t/2}\mathbf 1_{t\le0}$; its majorant is $e*\tilde e=e^{-\abs u/2}$.
Twisted by $y=\log x$ it gives the weights $\Lambda(n)\log(x/n)/\sqrt x$ on $n\le x$, i.e.\ $\Lambda(n)/\sqrt n$ times
$(n/x)^{1/2}\log(x/n)$. This is a classical Riesz-type weight (compare Bach \cite{bach1990} and \cite[Lemma~2.2]{lls2015}); it is the
member $a=\frac12$ of the family $\Lambda(n)(n/x)^a\log(x/n)$ in Bach's method as presented in \cite{mcgowntrevino2021}.
So the limiting constant $1$ is reached by a classical weight, and the windows of Proposition~\ref{prop:convsq} approach it as $L\to\infty$.
The inequality is a sharp point evaluation for functions whose Fourier transform extends to the upper half-plane. In the one-sided
least-non-residue problem, Carneiro, Milinovich, Quesada-Herrera and Ramos \cite{cmqr2024} find the limiting constant $C(\infty)=1$, with
the same extremal function, by a different (dual) argument. We do not claim novelty for Proposition~\ref{prop:barrier}. Its content for us is
that the constant $1$ is the exact limit of the explicit-formula method for prime-supported small-bias measures. This is a limit of
the method, not a lower bound for the true bias.

\begin{proof}[Proof of Corollary~\ref{cor:asymp}]
\emph{Part 1} is Proposition~\ref{prop:barrier}. \emph{Part 2.} A positive-definite kernel supported in $[-2a,2a]$ has $T=2a$, so
\eqref{eq:leading} and Proposition~\ref{prop:opt} give $c\ge e^a/(a+\sinh a)$. Minimizing:
$(\log c)'(a)=1-(1+\cosh a)/(a+\sinh a)$ vanishes iff $a+\sinh a=1+\cosh a$, i.e.\ $a-1=\cosh a-\sinh a=e^{-a}$. There
\[
  c(a)=\frac{e^a}{1+\cosh a}=\frac{2e^a}{(e^{a/2}+e^{-a/2})^2}=\frac2{(1+e^{-a})^2}=\frac2{a^2}.
\]
Since $c(a)\to\infty$ as $a\to0^+$ and $c(a)\to2$ as $a\to\infty$, while $c(a_*)<2$, the critical point $a_*$ is the global
minimum. \emph{Part 3.} Given $\delta>0$, choose $L$ with $1/(1-e^{-L/2})<1+\delta/2$ and a smooth $f\ge0$ supported in
$[-L/2,0]$, close enough to $e^{t/2}\mathbf 1_{[-L/2,0]}$ in $L^2$ that $c<1+\delta$ in Proposition~\ref{prop:convsq}. For this fixed
kernel every other term of the bound is $O_L(1)$ in the numerator, and $O_L(1)+O(\log q\log x/\sqrt x)$ in the denominator, which is
$o(\sqrt x)$ once $x/\log^2q\to\infty$.
\end{proof}

\subsection{General nonnegative kernels: numerics}
Only $K\ge0$ is needed for the weights to form a measure. Over nonnegative $K$ supported in a window of length $L$, the infimum of
the last expression in \eqref{eq:leading} (normalized by $e^{T/2}$) is a linear program after discretization
(\lean{kernel_lp.py}: piecewise-linear $K$, a polygonal model of $\abs{\cdot}$, solved with HiGHS; we report the exact ratio
of each LP solution). It lies strictly between the barrier and the convolution squares:
\begin{center}
\begin{tabular}{lccccc}
\toprule
$L$ & $4$ & $5.11$ & $6$ & $8$ & $12$\\
\midrule
positive-definite, $e^a/(a+\sinh a)$, $a=L/4$ & $1.250$ & $1.224$ & $1.235$ & $1.313$ & $1.543$\\
convolution squares, $1/(1-e^{-L/2})$ & $1.157$ & $1.084$ & $1.052$ & $1.019$ & $1.003$\\
nonnegative, LP (numerical) & $1.103$ & $1.047$ & $1.024$ & $1.006$ & $1.000$\\
\bottomrule
\end{tabular}
\end{center}
These kernels are not convolution squares, so for them the zero sum has to be bounded through zero counting, as in \cite{lls2015},
and not by positivity. We have not made that route explicit.

\begin{remark}
Lemma~\ref{lem:ranktwo} is formalized. Its specialization to $L^2[-a,a]$, the factorization theorem, Propositions~\ref{prop:convsq}
and~\ref{prop:barrier}, and the asymptotic statement are not. In the one-sided problem of \cite{cmqr2024} the sign conditions are
adapted to a single non-residue condition, and the optimum for the least quadratic non-residue is $0.7615\ldots\log^2p$. Our problem
bounds all characters simultaneously, in both directions.
\end{remark}

\section{Sampling, seeds and complexity}\label{sec:complexity}

This section is informal. Nothing in it is formalized beyond the bias bounds it invokes.

\paragraph{Computational model.}
Integers are given in binary. The modulus $q$ is given only as an integer: neither its factorization,
nor $\varphi(q)$, nor any generator of $\Gq$ is needed. The only use of $q$ is a coprimality test.

\paragraph{The sample space is $q$-independent.}
For fixed $(a,y)$ the weights $w_{a,y}(n)$ do not depend on $q$; $q$ enters only through the filter
$\gcd(n,q)=1$. Let $W=\{n\in[xe^{-2a},xe^{2a}]:\Lambda(n)>0\}$. A sieve up to $M=\lfloor e^{y+2a}\rfloor+1$
lists $W$ and the weights in $O(M\log\log M)$ arithmetic operations. Each weight needs one evaluation of
$k_a$, a smooth function of one real variable, computed once per point to any desired precision. For the
certificate of Corollary~\ref{cor:cert}, $M=166254$ and $\abs W$ is about $1.5\cdot10^4$.

\paragraph{Exact sampling and seed length.}
To output a sample with $s$ uniform random bits, round the normalized weights to multiples of $2^{-s}$ and
use inverse-CDF lookup in the precomputed table. If the rounding error in total variation is at most $\delta$,
the bias increases by at most $2\delta$. Taking $s=\lceil\log_2\abs W\rceil+\lceil\log_2(1/\delta)\rceil$ suffices.
With support $X\asymp\log^2q/\varepsilon^2$ (Corollary~\ref{cor:asymp}) this gives an $(\varepsilon+2\delta)$-biased
generator for $\Gq$ with seed length
\[
  s=2\log_2\log q+2\log_2(1/\varepsilon)+\log_2(1/\delta)+O(1)
\]
(the $O(1)$ absorbs the leading constant and the prime-power density). Each sample then costs one table lookup and one
gcd with $q$. The lower bound for any $\varepsilon$-biased set in a group of order $\varphi(q)$ is
$\log_2\log q+2\log_2(1/\varepsilon)-O(\log\log(1/\varepsilon))$ \cite{alonroichman1994,naornaor1993}. The
factor $2$ in front of $\log_2\log q$ is the $(\log q)^2$ barrier of Section~\ref{sec:limits}.

\paragraph{Non-units and $q$-obliviousness.}
Sampling from $W$ without the coprimality filter gives a measure that does not depend on $q$ at all. The
non-units of $W$ carry relative mass at most $U_q(a,y)/\mathrm{Den}$, which Theorem~\ref{thm:closed} bounds by
$U_B/\mathrm{den}_B$. For $q\le2^{64}$ with the parameters of Corollary~\ref{cor:cert} this is at most
$37.15/555.8<0.067$. So a single, $q$-oblivious table serves all moduli $q\le2^{64}$ simultaneously, at the
cost of a small failure probability.

\paragraph{Expander walks.}
By Corollary~\ref{cor:cert}, the random walk on $\Gq$ that multiplies by an independent sample at each step has
second eigenvalue at most $\frac12$ for every $q\le2^{64}$. The standard expander-walk Chernoff and hitting
bounds \cite{hlw2006} therefore apply to walks on $\Gq$ driven by these samples. A walk of length $t$ from a
uniform start uses $\log_2\varphi(q)+t\,s$ random bits instead of $(t+1)\log_2\varphi(q)$. Drawing the uniform
start does require sampling $\Gq$, e.g.\ by rejection with gcd tests.

\paragraph{What this is not.}
These are statistical guarantees against multiplicative characters of $\Gq$. They give no
indistinguishability against other tests (e.g.\ additive structure of $n$, or $n$ itself), and no
cryptographic security.

\section{Experiments}\label{sec:experiments}

All experiments are reproducible from \lean{research/grh-pseudorandomness/experiments/} in the accompanying repository. For a prime $q$, the group $\Gq$
is cyclic of order $q-1$. After fixing a primitive root and the discrete-logarithm table, the biases against
all $q-1$ characters are the entries of a single length-$(q-1)$ FFT. The biases are therefore computed
exactly, up to floating-point rounding, not estimated by sampling. These numbers describe typical moduli
and are not proofs.

\paragraph{Primes versus random sets.}
Table~\ref{tab:typical} compares the uniform measure on primes $p\le x$ with a uniformly random multiset of
residues of the same size (\lean{prime_bias.py}). On typical prime moduli up to $10^6$, primes are as good as
random sets, and in fact slightly better. The normalized bias $\lambda_*\sqrt{\abs{\mathrm{supp}}/\log q}$ is
$0.85$--$0.95$, the Alon--Roichman scale. The GRH-provable scale $\log q/\sqrt x$ is pessimistic by roughly
$\sqrt{\log q}$ on typical moduli; adversarial moduli at accessible sizes did not change this
(\lean{adversarial_q.py}).

\begin{table}[ht]
\centering
\begin{tabular}{rrrr}
\toprule
$q$ & $x$ & $\lambda_*$, primes $\le x$ & $\lambda_*$, random multiset\\
\midrule
10007 & 3200 & 0.1182 & 0.1474\\
100003 & 3200 & 0.1378 & 0.1619\\
1000003 & 3200 & 0.1487 & 0.1631\\
\bottomrule
\end{tabular}
\caption{Exact bias of primes $p\le x$ against all nontrivial characters, versus a random multiset of the same
size ($452$).}\label{tab:typical}
\end{table}

\paragraph{The bounds dominate the truth.}
For the theorem's own measure we compared the exact bias with the theorem's bound. In
\lean{cosh_kernel_check.py} the bound uses the explicit-formula constants, evaluated by quadrature: the
measured bias is $0.50$--$0.74$ of the bound in $36/36$ cases ($q\le3.8\cdot10^6$, $a\in\{1,1.5,2\}$). In
\lean{sharp_check.py} the bound is $\mathrm{Num}/\mathrm{Den}$ of Theorem~\ref{thm:sharp} with quadrature
constants: the ratio is $0.03$--$0.27$ in $8/8$ cases. A violation in either case would have revealed an
error in the numerics, since the theorem itself is machine-checked. For the certified measure of
Corollary~\ref{cor:cert}, the exact bias for $q\in\{10007,100003,1000003,4000037\}$ is
$0.021$, $0.026$, $0.032$, $0.032$ (\lean{certificate_check.py}), against the worst-case guarantee $\frac12$. For the
smaller measure of Corollary~\ref{cor:convcert} it is $0.037$, $0.047$, $0.048$, $0.055$ (\lean{convsq_certificate_check.py}).

\paragraph{What the theorems certify.}
Table~\ref{tab:params} gives the support needed to certify bias $\le\varepsilon$ for all $q\le2^{64}$.

\begin{table}[ht]
\centering
\begin{tabular}{lrrr}
\toprule
 & $\varepsilon=0.5$ & $\varepsilon=0.1$ & $\varepsilon=0.01$\\
\midrule
Theorem~\ref{thm:A} via Corollary~\ref{cor:A} (formal) & $9.7\cdot10^{15}$ & $3.5\cdot10^{17}$ & $5.6\cdot10^{19}$\\
Theorem~\ref{thm:closed}, closed form (formal once instantiated) & $1.3\cdot10^{5}$ & $1.8\cdot10^{6}$ & $1.5\cdot10^{8}$\\
Corollary~\ref{cor:cert} (formal, instantiated) & $1.66\cdot10^{5}$ & --- & ---\\
Theorem~\ref{thm:convclosed}, convolution square (formal once instantiated) & $4.2\cdot10^{4}$ & $6.7\cdot10^{5}$ & $5.9\cdot10^{7}$\\
Corollary~\ref{cor:convcert} (formal, instantiated) & $4.84\cdot10^{4}$ & --- & ---\\
Theorem~\ref{thm:sharp}, constants by quadrature & $3.5\cdot10^{4}$ & $5.5\cdot10^{5}$ & $4.7\cdot10^{7}$\\
Leading term $c^{*2}\log^2q/\varepsilon^2$, positive-definite (Cor.~\ref{cor:asymp}) & $1.2\cdot10^{4}$ & $2.9\cdot10^{5}$ & $2.9\cdot10^{7}$\\
Leading term $\log^2q/\varepsilon^2$, barrier (Cor.~\ref{cor:asymp}) & $7.9\cdot10^{3}$ & $2.0\cdot10^{5}$ & $2.0\cdot10^{7}$\\
\bottomrule
\end{tabular}
\caption{Support $X$ (largest integer in the support) sufficient for bias $\le\varepsilon$ for every $q\le2^{64}$.
Rows two and four give the smallest support for which the closed-form ratio is $\le\varepsilon$, found numerically
(\lean{closed_form_params.py}, \lean{convsq_formal_params.py}); their numerical inequalities are machine-checked only for the
instances in rows three and five, which use round parameters ($y=2\log150$; $y=2\log220$, $R=12$). The last two rows omit
lower-order terms.}\label{tab:params}
\end{table}

\section{The Lean formalization}\label{sec:formal}

\subsection{Environment and dependency}
The development lives in \lean{research/grh-pseudorandomness/lean/} of the accompanying repository. It uses Lean v4.35.0 (commit \lean{eb6cafb7}) and
Mathlib commit \lean{56a00585}, the exact pins of the compact edition of the GRH proof,
\lean{github.com/lsernau/grh-lean} at commit \lean{72fead10}, which is a path dependency \cite{grhlean}. All
ten package revisions agree with that edition's manifest. GRH enters in exactly two places:
\begin{itemize}
\item \lean{GRHPseudo.grh}, defined as \lean{LFunction_ne_zero_of_half_lt_re}, of type
  \lean{DirichletLNonvanishingRightOfHalf}
  from the GRH project's \lean{Statements.lean}, which imports only Mathlib;
\item \lean{GRHPseudo.Main.zerosOnLine}, defined as \lean{LFunction_eq_zero_re}, together with
  \lean{riemannHypothesis}.
\end{itemize}
Every other result is proved with these as explicit hypotheses, and the unconditional main theorems are
one-line applications of them. We rebuilt the compact edition from source: 9843 jobs, no errors.

\subsection{Architecture}
\begin{center}\small
\begin{tabular}{p{0.46\textwidth}p{0.46\textwidth}}
\toprule
Module & Content\\
\midrule
\lean{GRH/Interface}, \lean{GRH/Discharge} & hypothesis interface; discharge by the GRH proof\\
\lean{AnalyticNumberTheory/LGrowth} & $\abs{L(s,\chi)}\le1+4q\abs s$ on $\operatorname{Re}s\ge\frac12$\\
\lean{AnalyticNumberTheory/LogDerivGRH} & $L'/L$ bound right of $\frac12$ (uses GRH)\\
\lean{AnalyticNumberTheory/Contour}, \lean{PrimeCharSum} & line shift; Theorem~\ref{thm:A}\\
\lean{Sharp/Twist}, \lean{TwoSided}, \lean{MainTerm} & Lemma~\ref{lem:twosided}; main term from RH\\
\lean{Sharp/GammaBounds}, \lean{FiniteSums} & explicit gamma sides; finite prime sums\\
\lean{Sharp/Assembly} & Theorem~\ref{thm:sharp}\\
\lean{Sharp/KernelBounds}, \lean{ConstBounds}, \lean{Certificate} & Theorem~\ref{thm:closed}\\
\lean{Sharp/Certificate64} & Corollary~\ref{cor:cert}\\
\lean{Sharp/Majorant}, \lean{GammaSharp} & positivity with a majorant; sharper gamma bound\\
\lean{Sharp/ConvKernel}, \lean{ConvAssembly}, \lean{ConvCertificate} & Theorem~\ref{thm:convclosed}, Corollary~\ref{cor:convcert}\\
\lean{Pseudorandomness/*} & bias, Lemma~\ref{lem:eigen}, Lemma~\ref{lem:ranktwo}, Corollary~\ref{cor:A}\\
\lean{MainTheorems/Conditional}, \lean{Unconditional} & the statements of Section~\ref{sec:results}\\
\bottomrule
\end{tabular}
\end{center}
Our own code is about $4400$ lines in $27$ files. A further $2200$ lines are vendored, with attribution, from
\cite{grhmiller,schoenfeldlean,rhprimeerror} and the readable GRH edition.

\subsection{Porting the explicit formula}
The Weil explicit formula of \cite{zeta23} (120 modules, $43{,}000$ lines) targets Lean v4.33.0-rc2 and
Mathlib \lean{51e6992e}. We ported its closure to our pin. Only $12$ files needed changes, in $23$ hunks;
the complete diff against upstream is \lean{verification/zeta23-port.diff}. The changes fall into a few
kinds:
\begin{itemize}
\item a tactic file deleted from Mathlib (\lean{linear_combination'}), which we vendored from the old Mathlib;
\item lemma renames (\lean{Finset.prod_le_prod} $\to$ \lean{prod_le_prod}$_0$, \lean{prod_le_one} $\to$ \lean{prod_le_one}$_0$);
\item \lean{logDeriv_mul}/\lean{logDeriv_prod} applied to pointwise lambdas, which must now be folded first;
\item \lean{positivity} no longer seeing through \lean{id}, and one \lean{fun_prop} measurability goal;
\item \lean{field_simp} normal forms;
\item \lean{autoImplicit} enabled for this library, as upstream.
\end{itemize}
None of the changes alters a statement.

\subsection{Verification}
\begin{center}\small
\begin{tabular}{p{0.58\textwidth}p{0.34\textwidth}}
\toprule
Check & Result\\
\midrule
\lean{#print axioms}, 31 theorems including all of Section~\ref{sec:results} & \lean{propext}, \lean{Classical.choice}, \lean{Quot.sound}\\
all axioms in the environment (Mathlib, GRH proof, port, ours) & Lean core only\\
\lean{leanchecker} replay, 32 project modules & all pass\\
\lean{leanchecker} replay, 120 ported modules & all pass\\
source scan for \lean{sorry}, \lean{axiom}, \lean{native_decide} & none\\
\bottomrule
\end{tabular}
\end{center}
\lean{leanchecker} replays the declarations of each named module; imported modules, including the GRH proof,
are trusted through their \lean{.olean} files, which we built from source. Kernel replay of the GRH proof
itself, including Comparator, nanoda and lean4lean runs, was done by others
\cite{sutherlandgrh,sutherlandshortgrh,rubinsteinpage}. We have not yet run Comparator against a
Mathlib-only statement file for our theorems.

\subsection{Faithfulness of statements}
The definitions that occur in the main statements are short and appear in full in the source:
\lean{bias}, \lean{charSum}, \lean{mass} (\lean{Pseudorandomness/Bias}); \lean{sampleSpace}, \lean{wt},
\lean{cutoff}, \lean{kern}, \lean{testF} (\lean{Sharp/Assembly}); \lean{numB}, \lean{denB}, \lean{UB}, \lean{Kc},
\lean{Cc} (\lean{Sharp/Certificate}, \lean{ConstBounds}); and \lean{y150} (\lean{Sharp/Certificate64}). For the convolution-square results they are \lean{Conv.F}, \lean{Conv.f},
\lean{Conv.K}, \lean{Conv.wt}, \lean{Conv.sampleSpace}, \lean{Conv.cutoff} (\lean{Sharp/ConvKernel}, \lean{ConvAssembly}) and
\lean{Conv.G}, \lean{Conv.UB}, \lean{Conv.numB}, \lean{Conv.denB}, \lean{Conv.y220} (\lean{Sharp/ConvCertificate}). The test
function \lean{testFnN} is grh-miller's and is defined through Mathlib's \lean{Real.smoothTransition}. One
convention differs from the paper: the Lean sample space is $\{0\le n<M\}$, whose term $n=0$ has $\Lambda(0)=0$.

\section{Limitations and open problems}\label{sec:limits}

\paragraph{Cutoff, cardinality and seed length.}
Three sizes of a measure $\mu$ on prime powers must be kept apart: its \emph{cutoff} $X=\max\operatorname{supp}\mu$, its
\emph{cardinality} $\abs{\operatorname{supp}\mu}$, and the \emph{seed length} (or entropy) needed to sample it. Our theorems and the
barrier below concern the cutoff certified by the explicit-formula method. A measure with cutoff $X$ has cardinality at most about
$X/\log X$, but thinning or reweighting can make the cardinality much smaller while keeping the cutoff, and a small set can contain
very large primes.

\paragraph{The $(\log q)^2$ barrier.}
Every prime-supported construction analysed through the explicit formula has guaranteed bias
$\asymp\log q/\sqrt X$ in terms of its cutoff $X$. The reason is that about $(\log q)/2\pi$ zeros per unit height can conspire;
spreading the weight over scales changes this only by constants. A measure on primes with cutoff $(\log q)^{2-\delta}$ and bias
$<1$ for all $q$ and all $\chi\ne\chi_0$ would put a prime $\le(\log q)^{2-\delta}$ outside every proper subgroup of $\Gq$, since its
support cannot lie in the kernel of any nontrivial character. That would improve the GRH bound for the least prime outside a
subgroup, which is open \cite{lls2015,cmqr2024}. Montgomery's $\Omega(\log q\log\log q)$ \cite{montgomery1971} limits what is
possible even in principle. Empirically, typical moduli behave like random sets (Section~\ref{sec:experiments}): a random set of
cardinality $\asymp\log q\log X/\varepsilon^2$ is typically $\varepsilon$-biased.

\paragraph{What is formal and what is not.}
Theorems~\ref{thm:A}, \ref{thm:sharp}, \ref{thm:closed} and Corollaries~\ref{cor:A}, \ref{cor:cert} are formal.
Proposition~\ref{prop:opt} is formal only in its abstract core (Lemma~\ref{lem:ranktwo}). Propositions~\ref{prop:convsq}
and~\ref{prop:barrier}, Corollary~\ref{cor:asymp} and Section~\ref{sec:complexity} are informal. Of the convolution squares, only the
family $f_A$ of Theorem~\ref{thm:convclosed} is formalized. Only one numerical instance of Theorem~\ref{thm:closed} is
machine-checked. The closed form loses a factor of about $4$ in support relative to the quadrature constants
(Table~\ref{tab:params}), mostly through the bound on $\mathrm{SP}(a)$ and the decay constants in $I_a(\beta)$.

\paragraph{Trust.}
Our unconditional theorems inherit the trust base of the Lean GRH proof. That proof has passed extensive
machine checks, but at the time of writing no mathematician has read it through, and its companion paper
does not contain a proof of the main inequality \cite{rubinsteinpage}.

\paragraph{Open problems.}
\begin{enumerate}
\item Formalize Propositions~\ref{prop:opt}, \ref{prop:convsq} and~\ref{prop:barrier}, and certify further instances,
  e.g.\ $q\le2^{1024}$ or $\varepsilon=10^{-2}$.
\item Determine the optimal nonnegative kernel on a window of finite length $L$. At $L=4$ it strictly beats convolution squares,
  also among $C^\infty$ kernels and through positivity with a majorant (Propositions~\ref{prop:separation} and~\ref{prop:majorant};
  formal with the constant $\frac98$ for a piecewise-linear kernel, computer-assisted for the sharper constants and the $C^\infty$
  kernel). Certified lower bounds (for instance from the dual problem (EP2) of \cite{cmqr2024}) and the dependence on $L$
  remain open. An explicit bias theorem from such a kernel would need the positivity lemmas for a general majorant
  (Proposition~\ref{prop:majorant}), but no explicit zero counting.
\item Extend to other groups whose characters are Dirichlet characters, such as quotients of $\Gq$ and abelian
  Galois groups over $\mathbb Q$. For nonabelian groups, e.g.\ $\mathrm{Aff}(\Z/q\Z)$, combine with zig-zag or
  semidirect-product arguments.
\item The class-group and isogeny-graph expanders of \cite{jmv2009} need GRH for Hecke $L$-functions of
  imaginary quadratic fields, which remains open.
\end{enumerate}

\paragraph{Acknowledgements and attribution.}
GRH for Dirichlet $L$-functions in Lean is due to Alpöge, with Claude \cite{alpoge2026,grhlean}. The explicit
formula library is \cite{zeta23}; the positivity and gamma-side lemmas and the test function are from
\cite{grhmiller}; the Perron formula is from \cite{schoenfeldlean}; the Borel--Carath\'eodory log-derivative
lemma is from \cite{rhprimeerror}. All are Apache-2.0.

\paragraph{Use of AI.}
This work was carried out with Claude Opus 5.5 (Anthropic), used through Claude Code. Claude proposed and checked
candidate results, wrote the Lean formalization, the numerical experiments and drafts of the text, and searched the
literature. The author chose the problems, directed the work and reviewed the results, and takes responsibility for
all claims. The results marked formal are checked by the Lean kernel, with the axiom audit of
Section~\ref{sec:formal}. Informal results, the exposition and the comparison with the literature rest on the
author's review. Errors found during the work, including an incorrect optimality claim in an earlier draft, are
recorded in the text and in the repository history.


\bibliographystyle{amsplain}
\bibliography{references}

\end{document}
